\chapter{Experimental Setup}
\label{ch:experimental_setup}

\section{Dataset: Objaverse}
\label{sec:dataset_objaverse}

To gauge the structural alignment between 3D object shapes and text descriptions, a massive dataset of multimodal paired data is requisite. For this effort, we elected to utilize Objaverse \cite{deitke2023objaverse}, which is the most extensive public repository of multimodal 3D models paired with natural language captions generated from the Cap3D captioning framework \cite{luo2023scalable}. Unlike earlier, restricted synthetic datasets only supporting limited classes of fixed, rigid objects, Objaverse provides a comprehensive, realistic geometric and semantic distribution well-suited for analyzing the correlation of the various cross-modal latent spaces.

\begin{figure}[htpb]
    \centering
    \includegraphics[width=\linewidth]{objaverse.png}
    \caption{Sample objects from the Objaverse dataset alongside their textual annotations. This paired structure provides the foundational multimodal data required to supervise and evaluate 3D-text latent space alignment.}
    \label{fig:objaverse_sample}
\end{figure}

\subsection{Feature Extraction Pipeline}
\label{subsec:feature_pipeline}

The experimental workflow was instantiated on the University of Padova’s high-performance computing (HPC) cluster, \texttt{labsrv7}. We distributed feature extraction as well as pair alignment across NVIDIA A100 (20GB \gls{mig}), Tesla V100 and Tesla T4 GPUs via use of the Slurm workload management system.

We begin by sampling $100{,}000$ 3D objects from the Objaverse dataset. Before extracting geometric features, each 3D mesh was standardized: we uniformly sampled $2{,}048$ surface points, then centered each sample by subtracting the point-cloud’s mean and normalized by taking the maximum value along its euclidean dimensions. The feature extraction pipeline efficiently processes objects in batches of $10{,}000$ objects and features an auto-resume mechanism in case a node failure or job interrupt occurs. After filtering out meshes containing an error and samples with insufficient or nonsensical associated text, we arrived at a set of $46{,}675$ paired samples that had a clean textual annotation for all items across every modality.

\paragraph{Preprocessing Details.} Beyond the point-cloud standardisation described above, each modality underwent an additional normalisation step prior to any alignment method being fitted. All 3D and text feature vectors were $L_2$-normalised to unit norm immediately after extraction, ensuring that subsequent cosine-similarity-based retrieval metrics (Section~\ref{subsec:topk_acc}) are not confounded by differences in the absolute magnitude of embeddings produced by architecturally distinct encoders. For methods requiring zero-mean input --- specifically the raw-feature ablation of Test 12 in Chapter~\ref{ch:gw_alignment} --- a StandardScaler transformation (zero mean, unit variance per dimension, fitted on the anchor split only and applied unchanged to the query split) was additionally evaluated as an alternative to $L_2$-normalisation, to isolate whether the choice of normalisation scheme itself could account for observed performance differences. No outlier removal or clipping was applied to any feature vector, as the diagnostic purpose of this thesis requires that raw encoder outputs remain unmodified wherever possible.

\begin{figure}[htpb]
  \centering
  \includegraphics[width=1\textwidth]{fig_pipeline}
  \caption{Overview of the end-to-end feature extraction and alignment
  pipeline. Raw 3D objects and text captions are independently encoded
  by frozen uni-modal backbones into high-dimensional feature matrices.
  \gls{cca}, fitted on $n_A = 30{,}000$ anchor pairs, projects both
  modalities into a shared $k = 50$ dimensional subspace. The resulting
  representations are evaluated on 500 held-out query pairs using
  matching accuracy and Top-$K$ retrieval.}
  \label{fig:pipeline}
\end{figure}

\subsection{Anchor Pairs and Query Protocol}
\label{subsec:anchor_protocol}

In conjunction with the approach employed by \textcite{hadgi2025escaping}, we extract the alignment spaces and evaluate the representations themselves utilizing a strictly-defined subset partition of the dataset to define an anchoring subset and a querying subset, composed of $n_A = 30{,}000$ and $500$ matched pairs respectively. This allows explicit training for Canonical Correlation Analysis (CCA) projection and affine-mapping construction from the canonical directions whilst, at the same time, maintaining unseen objects used exclusively to evaluate the quality of the generated alignments. All evaluated quantitative results were produced by taking the mean over three independent data splits defined by these anchor/query partitions.

\paragraph{Hyperparameter Summary.} Table~\ref{tab:hyperparam_summary} consolidates the principal hyperparameters used across every alignment method evaluated in this thesis, most of which are introduced individually in their respective chapters but are collected here for ease of reference and reproducibility.

\begin{table}[htpb]
\centering
\renewcommand{\arraystretch}{1.15}
\resizebox{\textwidth}{!}{%
\begin{tabular}{llp{7.5cm}}
\toprule
\textbf{Method} & \textbf{Chapter} & \textbf{Key hyperparameters} \\
\midrule
CCA + Affine       & \ref{ch:baseline}   & $n_A=30{,}000$, $k=50$, 3 seeds, least-squares affine fit \\
KCCA               & \ref{ch:nonlinear}  & RBF kernel, $\gamma=0.1$, Nystr\"{o}m landmarks $m=500$, $n_A=30{,}000$ \\
RML                & \ref{ch:nonlinear}  & batch size 256, learning rate 0.01, 200 SGD iterations, Cholesky-factor parametrisation \\
SGW                & \ref{ch:gw_alignment} & 100 random projection directions; best of 288-configuration grid search over $k \in \{20,50,100\}$, $n \in \{1\text{k},5\text{k},30\text{k}\}$, $\lambda \in \{0,0.1,1.0,10.0\}$, 3 seeds (best: $k=20$, $n=5{,}000$) \\
RISGW              & \ref{ch:gw_alignment} & Stiefel manifold optimisation, 20 projections, 30 gradient iterations, $n=1{,}000$ \\
SS-RISGW           & \ref{ch:gw_alignment} & $\lambda \in \{0, 0.1, 1.0, 10.0\}$, $n=1{,}000$ anchor pairs \\
FGW                & \ref{ch:gw_alignment} & $\alpha=0.2$ (best), $k=50$ CCA subspace, $n_A=1{,}000$ \\
LR-GW              & \ref{ch:gw_alignment} & rank $r \in \{2,\ldots,200\}$ (best $r{=}50$), $\epsilon \in \{10^{-4},\ldots,5.0\}$, angular cost, OTT-JAX / LRSinkhorn solver \\
\bottomrule
\end{tabular}%
}
\caption{Consolidated hyperparameter summary across all alignment methods evaluated in this thesis. Full descriptions and ablation ranges are given in the corresponding chapter.}
\label{tab:hyperparam_summary}
\end{table}

\section{Encoders}
\label{sec:encoders_setup}

To thoroughly investigate how different neural architectures shape their resulting latent spaces, we set up a diverse array of feature extractors. Our selection strategy was not necessarily aimed at pushing state-of-the-art performance, but rather at understanding how architectural priors---such as unimodal versus multimodal pre-training---dictate the geometric topology of the embedded data.

\subsection{3D Encoders: PointNet and SparseConv}
\label{subsec:3d_encoders}

Processing 3D geometries introduces unique challenges, primarily because raw
point clouds lack a regular grid structure and are strictly permutation
invariant. To handle our standardised $2{,}048$-point samples, we deployed
two fundamentally different geometric backbones.

We relied on PointNet as our primary pre-trained extractor. PointNet directly
consumes the raw $(x, y, z)$ coordinates, passing them through shared
Multi-Layer Perceptrons (MLPs) and eventually using a symmetric max-pooling
function to aggregate global shape features. This process yields a robust
$1024$-dimensional feature vector. Its straightforward, global approach makes
it an excellent baseline for latent space analysis, as the resulting geometry
is not overly distorted by hierarchical inductive biases.

To contrast this, we also integrated a Sparse Convolutional Network
(SparseConv). Unlike PointNet's global pooling, SparseConv discretises the
3D space into a sparse voxel grid and applies convolutions exclusively to
occupied regions, which is theoretically much better at capturing fine-grained
local geometries. We configured this backbone to output $512$-dimensional
embeddings. Notably, as we will discuss in the baseline evaluations, we
deliberately utilised SparseConv without pre-training (random weight
initialisation).\footnote{Pre-trained SparseConv weights were requested from
the original authors of \textcite{hadgi2025escaping} but were not available
at the time of writing. The random-weight configuration therefore represents
a lower bound on SparseConv performance, and results are expected to improve
substantially once pre-trained weights become available.} This specific
experimental choice allows us to isolate and observe the raw alignment
capabilities of the spatial architecture itself, entirely separated from
learned feature representations.

\subsection{Text Encoders: OpenCLIP, RoBERTa, and BERT}
\label{subsec:text_encoders_setup}

For the textual modality, the central research question is whether purely
linguistic representations can be mapped to 3D geometries via optimal
transport, or if explicit visual pre-training is a strict prerequisite for
cross-modal alignment. To cover this spectrum, we extracted features using
three distinct transformer-based models.

Our strongest text backbone is the encoder derived from OpenCLIP
(specifically the \texttt{ViT-bigG-14} variant). Because \gls{clip} is
contrastively trained on massive image-text pairs, its latent space
inherently encodes physical and visual semantics. We hypothesise that this
underlying ``visual grounding'' makes its text embeddings significantly more
compatible with 3D structural spaces.

To evaluate purely linguistic spaces, we extracted text embeddings using \gls{bert}
and \gls{roberta}. These models were trained strictly on masked language modelling
objectives over massive text corpora, with absolutely no exposure to visual,
spatial, or geometric data. \gls{bert} provides our standard uni-modal baseline,
while \gls{roberta} represents a heavily optimised evolution of the same
architecture. By comparing the alignment behaviour of these pure language
models against \gls{clip}, we can accurately quantify how much of the cross-modal
translation success relies on the text encoder already ``knowing'' what the
physical world looks like --- a hypothesis that is directly validated by the
results presented in Chapter~\ref{ch:baseline}, where \gls{clip}-paired encoders
consistently and substantially outperform purely linguistic alternatives
across all evaluation metrics.

\section{Evaluation Metrics}
\label{sec:metrics}

Quantifying the success of cross-modal alignment requires robust metrics that capture both strict one-to-one correspondences and broader semantic retrieval capabilities. To ensure our evaluation remains strictly comparable with the supervised linear baseline, we assess the aligned representations using two primary criteria applied to our held-out query set of $500$ pairs.

\subsection{Matching Accuracy}
\label{subsec:matching_acc}

Matching accuracy is the most stringent metric we employ. We frame the evaluation as a bipartite matching problem: given a completely shuffled set of 500 query point clouds and their 500 corresponding text embeddings, the algorithm must reconstruct the exact original pairs. 

To achieve this, we first compute a dense pairwise distance matrix between all projected 3D and text features. We then apply the Hungarian algorithm (Linear Sum Assignment) to find the global permutation that minimizes the total cost across the entire batch. The matching accuracy is simply reported as the percentage of correctly recovered pairs. While mathematically elegant, this metric is exceptionally harsh. In a dataset as diverse as Objaverse, two distinct objects might share very similar geometries or semantic descriptions, meaning the algorithm is heavily penalized for confusing a generic "wooden chair" with another highly similar "wooden chair" in the batch.

\subsection{Top-K Retrieval Accuracy}
\label{subsec:topk_acc}

Because exact matching can be overly rigid for multimodal tasks, Top-$K$ retrieval serves as our primary indicator of practical alignment quality. Instead of forcing a global 1-to-1 assignment, retrieval tests the local neighborhood of each embedding. 

For a given 3D query shape, we rank all 500 text embeddings based on their cosine distance (or Euclidean distance, depending on the space topology) from the query point. If the ground-truth text caption falls within the top $K$ nearest neighbors, the retrieval is considered successful. We report results for $K \in \{1, 5, 10\}$. 

This metric is particularly insightful for analyzing latent spaces. A high Top-5 or Top-10 score, even with a low Top-1 score, suggests that the underlying geometry of the modalities is successfully clustering similar concepts together, even if the absolute nearest neighbor is slightly noisy. As we move into the Optimal Transport framework in the next chapters, these retrieval metrics will clearly expose whether Gromov-Wasserstein distances succeed at preserving these semantic neighborhoods or if they completely collapse the space.

\section{Diagnostic Metrics}
\label{sec:diagnostic_metrics}

While Matching Accuracy and Top-$K$ Retrieval (Section~\ref{sec:metrics}) quantify \emph{how well} an alignment method performs, they offer no insight into \emph{why} a method succeeds or fails. To support the theoretical and diagnostic analysis carried out throughout Chapter~\ref{ch:gw_alignment} --- in particular, the thirteen-test ablation study and the structural non-isometry argument of Section~\ref{sec:why_gw_fails} --- this section formally introduces two complementary diagnostic tools: Centered Kernel Alignment, used to directly quantify the structural similarity between two representation spaces, and transport plan entropy, used to characterise the internal structure of a learned coupling matrix independently of its downstream retrieval accuracy.

\subsection{Centered Kernel Alignment (CKA)}
\label{subsec:cka_definition}

Centered Kernel Alignment (CKA), introduced by \textcite{kornblith2019similarity}, measures the similarity between two sets of representations of the same $n$ samples, even when the representations have different dimensionalities. Given two feature matrices $X \in \mathbb{R}^{n \times d_1}$ and $Y \in \mathbb{R}^{n \times d_2}$, and two kernel matrices $K = k(X,X)$ and $L = l(Y,Y)$ computed over the two representation spaces, CKA is defined via the Hilbert-Schmidt Independence Criterion (HSIC) \cite{gretton2005measuring}:
\begin{equation}
    \mathrm{CKA}(K, L) = \frac{\mathrm{HSIC}(K, L)}{\sqrt{\mathrm{HSIC}(K, K) \, \mathrm{HSIC}(L, L)}}
    \label{eq:cka}
\end{equation}
where $\mathrm{HSIC}(K, L) = \frac{1}{(n-1)^2} \mathrm{tr}(K H L H)$, and $H = I_n - \frac{1}{n}\mathbf{1}\mathbf{1}^\top$ is the centering matrix. In the special, computationally convenient case of a linear kernel ($K = XX^\top$, $L = YY^\top$), Equation~\ref{eq:cka} reduces to the closed-form expression:
\begin{equation}
    \mathrm{CKA}_{\mathrm{linear}}(X, Y) = \frac{\| Y^\top X \|_F^2}{\| X^\top X \|_F \, \| Y^\top Y \|_F}
    \label{eq:linear_cka}
\end{equation}
which is the variant used throughout this thesis. CKA is bounded in $[0, 1]$, with a score of $0$ indicating no structural correlation between the two representation spaces, and a score of $1$ indicating that the pairwise geometric relationships of one space are perfectly recoverable from the other, up to an orthogonal transformation and isotropic scaling. Crucially, unlike the Gromov-Wasserstein distance of Section~\ref{sec:gw_distance}, computing CKA requires the two representation sets to be sampled over the \emph{same} underlying set of $n$ objects (i.e. it requires known correspondences), making it unsuitable as an unsupervised \emph{alignment} tool, but ideally suited as a supervised \emph{diagnostic}: because it is computed directly on the known 3D-text anchor pairs, it provides a ground-truth measurement of how structurally similar the two \emph{unaligned} latent spaces already are, independently of any alignment method applied to them. This is precisely the role CKA plays in the theoretical analysis of Section~\ref{subsec:isometry_bottleneck}, where the low CKA score reported by \textcite{hadgi2025escaping} between raw PointNet and CLIP embeddings ($\approx 0.12$) is used as independent corroborating evidence for the absence of structural isometry inferred from our own Gromov-Wasserstein ablation results.

\subsection{Transport Plan Entropy and Diagonal Concentration}
\label{subsec:entropy_definition}

A second diagnostic tool used extensively in the ablation study of Chapter~\ref{ch:gw_alignment} (Test 11, Section~\ref{subsec:library_analysis}) examines the internal structure of a learned transport plan $T \in \mathbb{R}^{n \times n}$ directly, independently of the retrieval accuracy it produces. Two complementary statistics are used.

\paragraph{Transport Plan Entropy.} As introduced in the entropic regularisation of Section~\ref{subsec:sinkhorn}, the (discrete) entropy of a transport plan is given by:
\begin{equation}
    H(T) = -\sum_{i,j} T_{ij} \log T_{ij}
    \label{eq:plan_entropy}
\end{equation}
A transport plan that concentrates all of its mass on a single, semantically correct permutation of matches (the ideal outcome for cross-modal retrieval) has minimal entropy: in the extreme case of a perfect permutation matrix, $H(T) = 0$. Conversely, a maximally uninformative transport plan --- one that spreads mass uniformly across every possible pairing, i.e. $T_{ij} = 1/n^2$ for all $i,j$ --- attains the theoretical maximum entropy $H_{\max} = \log(n^2)$. A learned transport plan whose entropy sits close to $H_{\max}$ is therefore statistically indistinguishable from random guessing, regardless of the retrieval accuracy computed from it: this is precisely the diagnostic signature used to characterise the failure of unsupervised Gromov-Wasserstein solvers in Test 11.

\paragraph{Diagonal-to-Uniform Concentration Ratio.} Because the ground-truth correspondence between a query set of $n$ 3D objects and their $n$ paired captions is, by construction, the identity permutation (query $i$ corresponds to caption $i$), a second, more interpretable diagnostic examines how much probability mass a transport plan places on its own diagonal relative to the uniform baseline:
\begin{equation}
    \rho(T) = \frac{\frac{1}{n}\sum_{i} T_{ii}}{\frac{1}{n^2}\sum_{i,j} T_{ij}}
    \label{eq:diagonal_ratio}
\end{equation}
which measures the mean diagonal mass of $T$ relative to what would be expected under a uniform, entirely uninformative coupling. A ratio of $\rho(T) = 1.0$ indicates a transport plan indistinguishable from uniform noise, while $\rho(T) > 1$ indicates a genuine (if imperfect) tendency to recover the correct correspondences. This ratio is reported alongside CCA+Affine's implicit transport structure and the unsupervised GW solvers' output plans in Section~\ref{subsec:library_analysis}, providing a direct, retrieval-metric-independent confirmation that unsupervised solvers fail to recover any meaningful structural signal.